% !TEX root = ../main.tex

\providecommand{\kBKM}{\kappa_{\mathrm{BKM}}}
\providecommand{\kch}{\kappa_{\mathrm{ch}}}
\providecommand{\kcat}{\kappa_{\mathrm{cat}}}
\providecommand{\kfib}{\kappa_{\mathrm{fiber}}}
\providecommand{\DBKM}{\Delta^{\mathrm{BKM}}}
\providecommand{\PhiPhys}{\Phi^{\mathrm{phys}}}
\providecommand{\AdS}{\mathrm{AdS}}
\providecommand{\Tr}{\mathrm{Tr}}

\section*{CHL dyon indices and Borcherds weights}
\label{wn:sec:chl-bps-bkm}

\paragraph{Primitive CHL geometry.}
Let
\[
  X^{(N)}=(K3\times E)/\mathbb{Z}_N,\qquad
  N\in\{1,2,3,4,6\},
\]
where \(g_N\) is a symplectic automorphism of the K3 factor and the
elliptic factor supplies the compatible free diagonal action. The
quotient is a smooth compact Calabi-Yau threefold in the primitive CHL
range. Let \(\DBKM_N\) be the Borcherds denominator attached to the
\(g_N\)-twined weak Jacobi input. The Borcherds input constants are
\[
  (c_1(0),c_2(0),c_3(0),c_4(0),c_6(0))
  =
  (10,8,6,4,2),
\]
hence Borcherds' weight theorem gives
\[
  \kBKM(\DBKM_N)=\frac{c_N(0)}2=(5,4,3,2,1).
\]
At \(N=1\), \(\DBKM_1=\Delta_5\), \(\mathrm{wt}(\Delta_5)=5\), and the
physical Igusa denominator is \(\Phi_{10}=\Delta_5^2\). For
\(N\in\{2,3,4,6\}\), the square-root Borcherds denominator \(\DBKM_N\)
and the square-normalized physical Siegel denominator \(\PhiPhys_N\)
are separate objects; the exact normalization comparison is an
additional datum.

\paragraph{Separate Borcherds catalogues.}
The constants above belong to the primitive CHL ladder indexed by
\(N\in\{1,2,3,4,6\}\). The Gritsenko-Clery diagonal-divisor, or dd,
catalogue is a separate eight-row paramodular atlas with its own row
indices, cover groups, multipliers, and weights. A comparison between a
CHL denominator and a Gritsenko-Clery dd row requires an explicit row
map and multiplier comparison; the tuple \((10,8,6,4,2)\) is the
primitive CHL tuple.

\paragraph{Mathieu twining census.}
The ordinary \(M_{24}\) twining table has \(26\) ATLAS conjugacy classes:
\[
\begin{gathered}
1A,\ 2A,\ 2B,\ 3A,\ 3B,\ 4A,\ 4B,\ 4C,\ 5A,\ 6A,\ 6B,\ 7A,\ 7B,\ 8A,\\
10A,\ 11A,\ 12A,\ 12B,\ 14A,\ 14B,\ 15A,\ 15B,\ 21A,\ 21B,\ 23A,\ 23B.
\end{gathered}
\]
The frame-shape rows most sensitive to misidentification are
\[
\begin{array}{c|c}
\text{class} & \text{frame shape} \\
\hline
3B & 3^8 \\
4A & 2^4 4^4 \\
4B & 1^4 2^2 4^4 \\
6B & 6^4 \\
10A & 2^2 10^2 \\
12A & 2^1 4^1 6^1 12^1 \\
12B & 12^2 .
\end{array}
\]
Each row satisfies \(\sum_i i a_i=24\) for frame shape
\(\prod_i i^{a_i}\). The primitive CHL ladder uses the chosen CHL
classes and their orbifold Jacobi inputs; the ordinary \(M_{24}\)
census remains independent.

\paragraph{Protected dyon index.}
For charges \((Q,P)\) in the CHL electric-magnetic Narain lattice, put
\[
  \Delta_N(Q,P)=Q^2P^2-(Q\cdot P)^2 .
\]
In a contour chamber \(\mathcal C\) avoiding walls of marginal
stability, the Jatkar-Sen quarter-BPS helicity supertrace is extracted
by
\[
d_N(Q,P;\mathcal C)
=
(-1)^{Q\cdot P+1}
\oint_{\mathcal C}
\frac{
e^{-\pi i(Q^2 z_1+2(Q\cdot P)z_2+P^2 z_3)}
}{
\PhiPhys_N(Z)
}
d^3Z,
\qquad
Z=\begin{pmatrix}z_1&z_2\\ z_2&z_3\end{pmatrix}\in\mathfrak H_2 .
\label{wn:eq:chl-dyon-count}
\]
The output is a protected signed index in the specified CHL chamber. At
\(N=1\) this is the Dijkgraaf-Verlinde-Verlinde contour formula for
\(1/\Phi_{10}=1/\Delta_5^2\). For \(N\in\{2,3,4,6\}\), the theorem-level
input is the CHL protected-index formula for the physical denominator
\(\PhiPhys_N\); comparison with the BKM square root is an additional
normalization theorem.

\paragraph{Large-charge saddle.}
The attractor-chamber saddle gives
\[
  \log d_N(Q,P;\mathcal C)
  =
  \pi\sqrt{\Delta_N(Q,P)/N}+O(\log \Delta_N).
\label{wn:eq:chl-leading-entropy}
\]
In the BKM square-root normalization, the automorphic weight
contribution to the logarithmic term is
\[
  \kBKM(\DBKM_N)-\frac32
  =
  \frac{c_N(0)}2-\frac32 .
\label{wn:eq:chl-bkm-log-weight}
\]
Conversion from \eqref{wn:eq:chl-bkm-log-weight} to the full quantum
entropy logarithmic coefficient requires the square-root comparison,
the contour measure, the Gaussian determinant, and the massless
zero-mode subtraction in one convention.

\paragraph{AdS and index scope.}
The holographic consequence supported by the protected dyon formula is
the leading entropy identity
\[
  S_{\mathrm{BH}}^{(N)}
  =
  \pi\sqrt{\Delta_N(Q,P)/N}.
\]
The Brown-Henneaux central charge is determined by the near-horizon
radius and the three-dimensional Newton constant of the chosen
D1-D5-KK-P host. The cohomological trace
\[
  \chi_{g_N}(K3)=\Tr(g_N\mid H^*(K3))
  \in\{24,8,6,4,2\}
\]
is the protected K3 fibre index entering the twined elliptic genus and
\(\kfib^{\mathrm{tw}}\).

\paragraph{Twined elliptic genus.}
The protected seed index is
\[
Z^{(g_N)}_{\mathrm{ell}}(K3)(\tau,z)
=
\Tr_{\mathrm{RR}}\!\left(
g_N q^{L_0-c/24}\bar q^{\bar L_0-\bar c/24}y^{J_0}
\right).
\label{wn:eq:chl-twined-elliptic-genus}
\]
It is a weak Jacobi form of weight \(0\) and index \(1\), with the
multiplier determined by the CHL congruence subgroup. A theorem-level
bulk comparison concerns this protected index, the contour chamber, and
the specified CHL host geometry; equality of ordinary partition
functions requires a localization theorem with matching boundary
conditions and measures.

\paragraph{Four invariants.}
The CHL quotient carries four separate invariants:
\[
\bigl(
  \kcat(X^{(N)}),
  \kch(\mathcal A_{X^{(N)}}),
  \kBKM(\DBKM_N),
  \kfib^{\mathrm{tw}}(X^{(N)})
\bigr)
=
\bigl(
  0,\ 0,\ c_N(0)/2,\ \chi_{g_N}(K3)
\bigr).
\label{wn:eq:chl-four-invariants}
\]
The equality \(\kcat(X^{(N)})=\chi(\mathcal O_{X^{(N)}})=0\) follows
from the finite etale cover \(K3\times E\to X^{(N)}\) and
\(\chi(\mathcal O_{K3\times E})=\chi(\mathcal O_{K3})\chi(\mathcal O_E)
=2\cdot0=0\). The compact chiral supertrace is
\[
  \kch(\mathcal A_{X^{(N)}})=
  \sum_q(-1)^q h^{0,q}(X^{(N)})=1-1=0
\]
for the compact Calabi-Yau threefold output. The Borcherds invariant is
the automorphic weight \(c_N(0)/2\). The fibre invariant is the Lefschetz
trace of \(g_N\) on K3 cohomology.

At \(N=1\), the compact \(K3\times E\) object also has the
four-construction menu
\[
  \{\kcat,\ \kappa_{\mathrm{ch}}^{\mathrm{Heis}},\ \kBKM,\ \kfib\}
  =
  \{0,3,5,24\}.
\]
Here \(\kch(K3\times E)=0\) is the compact Hodge supertrace, while
\(\kappa_{\mathrm{ch}}^{\mathrm{Heis}}=3\) is the Heisenberg-Mukai
specialization. The number \(5\) is the weight of \(\Delta_5\), and
\(24\) is the K3 fibre cohomological trace.

\paragraph{Operadic scope.}
Since \(\dim_{\mathbb C}X^{(N)}=3\), the specialized chiral output
\(\mathcal A_{X^{(N)}}\) is native \(E_1\)-chiral. Any \(E_2\) braided
structure belongs to the derived chiral centre
\[
  \mathcal Z(\operatorname{Rep}^{E_1}(\mathcal A_{X^{(N)}}))
\]
or to a completed Hall-Drinfeld double after the pairing, coproduct,
completion, and centre-continuity data are supplied.

\paragraph{Theorem boundary.}
The following data are required before a quantum-gravity or BPS-algebra
claim can be used as theorem-level input:
\begin{enumerate}
\item the precise protected index, including sign convention and charge
  lattice;
\item the contour or stability chamber;
\item the host geometry and quotient convention;
\item the comparison theorem relating the physical denominator
  \(\PhiPhys_N\) to \(\DBKM_N\);
\item the compact Hall-Drinfeld or BPS quantum-group construction,
  including orientation, primitive projection, PBW compatibility, and
  radical quotient.
\end{enumerate}
Consequently, the theorem-level content consists of the constants,
invariant separation, protected-index scope, and listed proof
obligations.

\paragraph{Sources.}
Borcherds 1998, Theorem 13.3; Gritsenko-Nikulin 1995; Gritsenko 1999;
Dijkgraaf-Verlinde-Verlinde 1997; Jatkar-Sen 2005; Sen 2007 and 2010;
Eguchi-Ooguri-Tachikawa 2011; the ATLAS \(M_{24}\) table;
Conway-Sloane; Gaberdiel-Hohenegger-Volpato 2012.
